\subsection{A common source-population scalar and quantum detector packet}
\label{sec:source-common-scalar-packet}

The prepared golden source addresses admit a local massive scalar action
whose field edges fit inside the causal construction's existing read radius.
The action below uses adjacent sorted source coordinates, rather than the
radial kernel. It therefore has its own coefficients and boundary condition.
The source addresses are protected preparation data, distinct from the
mutable field registers. No invariance of an address under arbitrary live
port repair is assumed.

Put $L=2/\sqrt{\phi+2}$ and use coordinates $y=x/L$, model time
$s=ct/L$, $\hbar=1$ and dimensionless mass one. For $q=F_n$, sort
$\xi_b=b\phi-\lfloor b\phi\rfloor$, $0\le b<q$, as
$0=x_0<x_1<\cdots<x_{q-1}<1$, and append $x_q=1$ as a fixed Dirichlet
boundary marker. This marker is also a source-coded point. The mutable
interior population $\{x_1,\ldots,x_{q-1}\}^3$ is contained in the original
golden product population. Each coordinate is the normalized readback of the
integer control $a=-\lfloor b\phi\rfloor,b$; the nodes are not replaced by
their nearby uniform-grid labels.

Write $h_i=x_{i+1}-x_i$ and $w_i=(h_{i-1}+h_i)/2$ for interior nodes.
Let $W=\operatorname{diag}(w_i)$ and let $K$ be the Dirichlet stiffness
matrix with diagonal $h_{i-1}^{-1}+h_i^{-1}$ and off-diagonal
$K_{i,i+1}=-h_i^{-1}$. On the tensor population use product masses and
the action
\begin{equation}
 \mathcal S_q=\frac12\int\left[
 \|\dot u\|_{W^{\otimes3}}^2-
 \langle u,(I+\mathsf L_1+\mathsf L_2+\mathsf L_3)u\rangle_{W^{\otimes3}}
 \right]ds,\qquad \mathsf L=W^{-1}K.
 \label{eq:common-scalar-action}
\end{equation}
Here $\mathsf L_d$ acts on coordinate $d$. Each varying field read is between
adjacent source sites; fixed boundary values contribute on-site terms.
The golden gap is at most $2/q$, so $h_i<q^{-1/2}$ for $q>4$.
The four finite certificates also check the stronger exact gap inequality
$qh_i^2<1$ directly. Thus these field edges are permitted by the radius
$a_q=L/\sqrt q$. This inclusion is not a construction of authenticated
time-step events. Subcycling and the custody of a count-derived clock
require an additional execution comparison.

Mass whitening gives the strictly positive real matrix
\[
 A_q=I+B\otimes I\otimes I+I\otimes B\otimes I+
 I\otimes I\otimes B,\qquad B=W^{-1/2}KW^{-1/2}.
\]
Canonical quantization is the oscillator Fock space with vacuum-subtracted
Hamiltonian $d\Gamma(\omega_q)$, $\omega_q=A_q^{1/2}\ge I$.
The comparison theory is the massive Dirichlet free scalar on the same
unit cube, with $\omega=(I-\Delta_D)^{1/2}$. Its nine comparison modes are
\[
 e_k(y)=\sqrt8\sin(k\pi y_1)\sin(\pi y_2)\sin(\pi y_3),
 \qquad \omega_k=\sqrt{1+\pi^2(k^2+2)},\quad 1\le k\le9.
\]
The remaining continuum and lattice modes belong to their respective full
Fock spaces; they are not discarded from an evolution residual.

\paragraph{A full frequency residual from small tridiagonal solves.}
Let $S$ sample these modes with square roots of the tensor masses,
$G=S^*S$, and suppose $\|G-I\|\le d<1$. Then
$j=SG^{-1/2}$ is an isometry. Put
\[
 c_d=(1-d)^{-1/2}-1,
 \qquad s_d=\sqrt{1+d}\,c_d,
 \qquad \Omega=\operatorname{diag}(\omega_k).
\]
For $f_k(x)=\sqrt2\sin(k\pi x)$ define
$s_k=W^{1/2}f_k|_S$ and
$r_k=W^{1/2}(\mathsf Lf_k-(k\pi)^2f_k)|_S$.
For each $k$ compute the two nonnegative quadratic forms
\[
 Q_{kj}=r_j^*[B+(1+\omega_k^2)I]^{-1}r_j,
 \qquad j\in\{k,1\},
\]
and set
\[
 b_k=\sqrt{Q_{kk}}\|s_1\|^2+
 2\sqrt{Q_{k1}}\|s_k\|\|s_1\|,\qquad
 r=\Big(\sum_k b_k^2\Big)^{1/2}.
\]
The quadratic forms require only the positive tridiagonal solve
$(K+(1+\omega_k^2)W)y=W(\mathsf Lf_j-(j\pi)^2f_j)$.

\begin{proposition}[Full lattice leakage and coefficient-sensitive control]
\label{prop:common-scalar-full-leakage}
With the definitions above,
\begin{align}
 \|\omega_qj-j\Omega\|&\le r(1-d)^{-1/2}+2\omega_9s_d, \notag\\
 \|(\omega_qj-j\Omega)v\|&\le\kappa(v):=
 \sum_k b_k|v_k|+(rc_d+2\omega_9s_d)\|v\|.
 \label{eq:common-scalar-kappa}
\end{align}
The bound satisfies $\kappa(Dv)\le\kappa(v)$ whenever $D$ is diagonal and
$|D_{kk}|\le1$. For $m=1$,
\begin{equation}
 \|\omega_q^{-1/2}jv-j\Omega^{-1/2}v\|
 \le\tfrac12\kappa(v).
 \label{eq:common-scalar-negative-power}
\end{equation}
\end{proposition}
\begin{proof}
For a positive scalar $\lambda$,
$(\sqrt\lambda-\omega_k)^2\le
(\lambda-\omega_k^2)^2/(\lambda+\omega_k^2)$.
Functional calculus applies this inequality to the full matrix $A_q$.
The stiffness residual on $S e_k=s_k\otimes s_1\otimes s_1$ is the sum
of its three tensor residuals. Minkowski's inequality in the
$(A_q+\omega_k^2)^{-1}$ norm, followed by
$A_q+\omega_k^2\ge B_d+(1+\omega_k^2)I$, bounds the full frequency
residual column by $b_k$. This step includes its component perpendicular
to the sampled nine-mode space. Hence
$\|\omega_qS-S\Omega\|\le r$ and its action on $v$ is bounded by
$\sum_kb_k|v_k|$.

Write $G^{-1/2}=I+C$, with $\|C\|\le c_d$ and $\|SC\|\le s_d$.
Expanding
\[
 \omega_qj-j\Omega=(\omega_qS-S\Omega)(I+C)+S[\Omega,C]
\]
gives both bounds. In the commutator estimate use
$\|S\|\|C\|\le s_d$; no commutation of $G$ and $\Omega$ is assumed.
Diagonal contraction decreases every term in $\kappa$.
Finally, the resolvent formula for the inverse square root gives
\[
 -\frac1\pi\int_0^\infty t^{-1/2}
 (\omega_q+t)^{-1}(\omega_qj-j\Omega)(\Omega+t)^{-1}v\,dt.
\]
Its norm is at most
$\kappa(v)\pi^{-1}\int_0^\infty t^{-1/2}(1+t)^{-2}dt
=\kappa(v)/2$.
\end{proof}

The isometry on the nine-mode comparison Fock space is $\Gamma(j)$.
On its sector with at most $N$ particles, the full generator residual is
at most $N[r(1-d)^{-1/2}+2\omega_9s_d]$, by the sum of the $N$ one-particle
residuals on symmetric tensors. This is an actual generator estimate,
not a bound on a compressed matrix. For the particular coherent states
below, using their coefficients directly gives a substantially smaller
comparison without an occupation cutoff.

\begin{proposition}[The same action has a free-field detector limit]
\label{prop:common-scalar-continuum-limit}
Along the golden refinements, the sampled isometries on every fixed finite
Dirichlet sine band have vanishing full frequency residual. For any fixed
finite time interval, the classical smeared responses to bounded,
piecewise-continuous $L^2$ momentum preparations, and the coherent-state
expectations of finitely many bounded smeared Weyl effects, converge to
those of the massive Dirichlet scalar on the cube. Smearings and their
products must be Riemann integrable; a finite collection of slab-boundary
jumps satisfies this condition. Initial position is zero for the momentum
preparations in this statement.
\end{proposition}
\begin{proof}
Let $\ell=\max_i h_i\to0$. Taylor's formula at an interior node gives
\[
 |\mathsf Lf(x_i)+f''(x_i)|\le
 \frac{h_{i-1}^2+h_i^2}{3(h_{i-1}+h_i)}\|f'''\|_\infty
 \le\frac\ell3\|f'''\|_\infty.
\]
This includes the generally nonzero error from unequal adjacent gaps.
The product trapezoidal rule has error at most
$\ell^2\sum_d\|\partial_d^2F\|_\infty/12$ on the unit cube.
Boundary masses may be included: the Dirichlet sine products vanish there,
so their contributions to the Gram matrix are zero.
Consequently $G-I=O(\ell^2)$ on every fixed band and the full stiffness
residual is $O(\ell)$. The square-root resolvent bound, or directly
$(\sqrt\lambda+\omega_k)^{-1}\le1/2$, makes the frequency residual
$O(\ell)$ as well. These are actual residual estimates for this action;
no ratio $H/a_q^2\to0$ is required.

Finite-band Duhamel and the negative-power estimate give uniform finite-time
convergence of $\omega_q^{-1}\sin(s\omega_q)$ between sampled finite
sine combinations, as well as the corresponding classical velocity
readouts. The graph kinetic energy norm is the same product-mass norm.
For general stated smearings, first approximate them in $L^2$ by finite
sine combinations. Their squared sampling errors converge to the squared
$L^2$ errors by product Riemann quadrature. The uniform bounds
$\|\omega_q^{-1/2}\|,\|\omega^{-1/2}\|\le1$ and unitarity control
the tails independently of $q$ and of time in the fixed interval. Take
$q\to\infty$ for a fixed sine approximation, then let its error tend to
zero. This also applies to the compact slab functions below.

The same argument controls coherent amplitudes, Weyl parameters, their
inner products and their squared norms. The Gaussian characteristic
function therefore gives the asserted bounded quantum readout convergence;
finite Weyl products use the Weyl relations and the same controlled inner
products. The limiting Fock theory is built on the full continuum
$L^2$ one-particle space. Each finite-band isometry has its own domain;
there is no isometry of that entire infinite one-particle space into a
single finite lattice space. This observable convergence does not assert
an interacting continuum or an operator-norm limit of all local effects.
\end{proof}

\paragraph{Compact preparation and a separated bounded detector.}
Normalize $f$ to unit $L^2$ norm, where
\[
 f(y)=C\sin(\pi y_1)(1+\cos(\pi y_1))^8
                 \sin(\pi y_2)\sin(\pi y_3),\qquad
 g(y)=f(1-y_1,y_2,y_3).
\]
Both have exactly the nine modes above; denote their coefficient vectors
by $a$ and $\widetilde a_k=(-1)^{k+1}a_k$.
The actual smearings are
\[
 f_c=f\,\mathbf1_{y_1\le9/20},\qquad
 g_c=g\,\mathbf1_{y_1\ge11/20}.
\]
Their supports are separated by $1/10$ in the normalized source coordinate,
or $L/10$ in the source metric.
Prepare $\exp(4i\Phi(f_c))|0\rangle$ and measure
$E_c=(I+\sin\Phi(g_c))/2$, with $0\le E_c\le I$.
Use the same sampled compact smearings and the same definitions in the
finite graph oscillator theory. These are field instruments with spatial
support, not a rank-one projector onto a named mode.
The unperturbed vacuum probability and the prepared probability at time
zero are both $1/2$.

For the smooth comparison, put
$b=\|\Omega^{-1/2}a\|/\sqrt2$ and $\alpha=4b$.
Its exact induced probability is
\begin{equation}
 P(s)-\tfrac12=
 \tfrac12e^{-b^2/2}\sin\!\left(
 4\sum_{k=1}^9(-1)^{k+1}a_k^2
                      \frac{\sin(\omega_ks)}{\omega_k}\right).
 \label{eq:common-scalar-probability}
\end{equation}
This follows from the coherent-state characteristic function. For two
coherent amplitudes differing by at most $d_\alpha$ and two Weyl parameters
differing by at most $d_b$, the corresponding effects obey the explicit
probability bound
\begin{equation}
 \mathcal E(b,\alpha,d_b,d_\alpha)=
 \tfrac14d_b(2b+d_b)+d_b(\alpha+d_\alpha)+bd_\alpha.
 \label{eq:common-scalar-gaussian-error}
\end{equation}
Indeed, $|e^{-x}-e^{-y}|\le|x-y|$ for $x,y\ge0$ and
$|\sin x-\sin y|\le|x-y|$ control the vacuum factor and the coherent
phase. All norms refer to the full one-particle Hilbert spaces after the
isometric comparison.

For the graph versus continuum smooth comparison, take
\[
 d_b=\frac{\kappa(a)/2+s_d}{\sqrt2},\qquad
 d_\alpha=4d_b+\frac{4|s|}{\sqrt2}\kappa(\Omega^{-1/2}a).
\]
Equation~\eqref{eq:common-scalar-negative-power} controls the initial
smearing maps. Duhamel's formula on the finite spectral comparison space
controls evolution; $\kappa(e^{-is\Omega}v)=\kappa(v)$ makes the bound
uniform in time. The graph unitary acts on every lattice mode.

Each compact-continuum smearing differs from its smooth version by squared
$L^2$ norm
$\theta=\int_{9/20}^1|\sum_k a_k\sqrt2\sin(k\pi x)|^2dx$.
Thus its extra $d_b$ is at most $\sqrt{\theta/2}$ and its extra
$d_\alpha$ at most $4\sqrt{\theta/2}$, using $\omega\ge1$.
The graph tails are computed separately from their actual lumped masses.
Their two values need not equal the continuum tail or each other.
Apply \eqref{eq:common-scalar-gaussian-error} to both localization steps
and add their errors to the smooth comparison error.
The continuum compact preparation has finite expected energy at most $8$
above the vacuum. The graph energy is $8\|f_c|_S\|_{W^{\otimes3}}^2$,
which the certificate computes separately.
A discontinuity at a slab boundary does not imply membership
in the Hamiltonian's operator domain. Only its form domain and unitary
norm preservation are used here; the Duhamel calculation uses the smooth
finite-mode vector. For any $E>0$, the continuum energy tail is at most $8/E$ by
the spectral Markov inequality.

\begin{theorem}[Resolved compact scalar quantum response]
\label{thm:source-common-scalar-response}
For the declared action, instruments and units above, the $q=233$ golden
population satisfies, uniformly for $19/20\le s\le1$,
\[
 |P_{233,c}(s)-P_c(s)|<0.050055,
 \qquad |P_{233,c}(s)-\tfrac12|>0.137754.
\]
The comparison error is smaller than the induced response throughout a
nonzero time interval. The response is to a spatially nonuniform compact
preparation at a spatially separated compact detector.
\end{theorem}
\begin{proof}
The executable certificate uses outward rational arithmetic, the exact
golden coordinates, positive tridiagonal elimination and the bounds above.
It encloses the continuum tail by $0.000222611$ and the full nine-mode
frequency residual by $0.053232$. The coefficient-sensitive preparation
residual is below $0.005360$. One hundred closed time cells cover the
entire interval; Taylor remainders enclose all trigonometric and exponential
values. The final inequalities include both sampled localization tails
and the continuum localization error. An independent verifier assembles
edge fluxes, solves the positive systems in reverse order, and recomputes
the probability bounds. The smaller populations $q=13,34,89$ are reported
without requiring them to satisfy the signal criterion.
\end{proof}

The tensor calculation uses one-dimensional arrays of length $232$, while
the mathematical finite model has $232^3$ oscillators and an
infinite-dimensional Fock space. No full read/write trace at that population
is claimed. The theorem is a free massive scalar comparison with prescribed
axes, boundary, action, protected addresses and model time. It does not
select these inputs from native repair, identify physical clock units,
prove an interacting or charged continuum, or identify Hamiltonian
substeps with the retained causal-count history.
